\documentclass[11pt,a4paper]{article}

\usepackage[utf8]{inputenc}
\usepackage[T1]{fontenc}
\usepackage[italian]{babel}
\usepackage{geometry}
\usepackage{graphicx}
\usepackage{amsmath}
\usepackage{amssymb}
\usepackage{booktabs}
\usepackage{tabularx}
\usepackage{array}
\usepackage{xcolor}
\usepackage{listings}
\usepackage{hyperref}
\usepackage{tikz}
\usepackage{caption}
\usepackage{ragged2e}
\usepackage{enumitem}
\usepackage{seqsplit}

\usetikzlibrary{arrows.meta,positioning,shapes.geometric,calc}

\geometry{margin=2.25cm}
\setlength{\parskip}{0.55em}
\setlength{\parindent}{0pt}
\setlength{\tabcolsep}{4.5pt}
\renewcommand{\arraystretch}{1.22}
\emergencystretch=3em
\sloppy

\captionsetup{font=small,labelfont=bf}

\hypersetup{
    colorlinks=true,
    linkcolor=blue!60!black,
    urlcolor=blue!60!black,
    citecolor=blue!60!black
}

\newcolumntype{Y}{>{\RaggedRight\arraybackslash}X}
\newcolumntype{L}[1]{>{\RaggedRight\arraybackslash}p{#1}}

\newcommand{\code}[1]{\texttt{\detokenize{#1}}}
\newcommand{\codeb}[1]{\texttt{\seqsplit{#1}}}
\newcommand{\cls}[1]{\code{#1}}
\newcommand{\clsb}[1]{\codeb{#1}}
\newcommand{\meth}[1]{\code{#1}}
\newcommand{\pkg}[1]{\code{#1}}
\newcommand{\formal}[1]{\emph{Formalizzazione}, #1}

\lstset{
    basicstyle=\ttfamily\footnotesize,
    breaklines=true,
    columns=fullflexible,
    frame=single,
    rulecolor=\color{black!25},
    backgroundcolor=\color{black!2},
    keepspaces=true
}

\tikzset{
    box/.style={
        draw,
        rounded corners=2pt,
        thick,
        align=center,
        text width=2.65cm,
        minimum height=0.82cm,
        inner sep=4pt,
        fill=blue!5,
        font=\scriptsize
    },
    widebox/.style={
        box,
        text width=3.15cm
    },
    data/.style={
        box,
        fill=green!8
    },
    policy/.style={
        box,
        fill=orange!10
    },
    result/.style={
        box,
        fill=purple!8
    },
    decision/.style={
        diamond,
        draw,
        thick,
        aspect=2.2,
        align=center,
        inner sep=1.5pt,
        fill=yellow!14,
        font=\scriptsize
    },
    arrow/.style={
        -{Latex[length=2.2mm]},
        thick
    },
    dashedarrow/.style={
        -{Latex[length=2.2mm]},
        thick,
        dashed
    }
}

\title{\textbf{Guida tecnica al sistema MA-GA}\\
\large Collegamento tra classi, metodi e formalizzazione}
\author{Repository \code{maga-core}}
\date{\today}

\begin{document}

\maketitle

\begin{abstract}
Questo documento spiega come il codice realizza il Mobility-Aware Genetic
Algorithm per il computation offloading nel continuum veicolo, edge e cloud.
Il focus e' pratico: cosa fa ogni classe, con chi comunica, quale decisione
prende e come questa decisione si collega alla formalizzazione. Il testo tiene
conto delle correzioni recenti, in particolare del provider filtrato e delle
nuove policy per quota di offloading, CPU e banda.
\end{abstract}

\tableofcontents
\newpage

\section{Scopo}

Il sistema prende uno snapshot del mondo simulato. Dentro lo snapshot ci sono
veicoli, task e candidati di esecuzione. Il GA costruisce una strategia:
per ogni task sceglie dove eseguirlo, quanta parte mandare in remoto e quante
risorse assegnare.

La formalizzazione descrive questo problema con variabili, vincoli e funzione
obiettivo. Il codice lo traduce in oggetti Java. La scelta finale non viene
presa da una sola classe. Nasce da una catena:

\begin{itemize}[leftmargin=*,nosep]
    \item lo snapshot descrive lo stato;
    \item il prefilter riduce i candidati poco utili;
    \item le policy generano scelte iniziali piu' sensate;
    \item il GA esplora lo spazio delle soluzioni;
    \item la fitness misura il costo;
    \item il repair corregge soluzioni non coerenti;
    \item il gestore temporale decide quando rieseguire tutto.
\end{itemize}

\section{Mappa generale}

La figura mostra il flusso principale. E' volutamente compatta. Serve a vedere
chi chiama chi, non ogni dettaglio interno.

\begin{figure}[h]
\centering
\resizebox{0.98\linewidth}{!}{%
\begin{tikzpicture}[node distance=1.15cm and 1.0cm]
    \node[data] (json) {file JSON\\snapshot + eventi};
    \node[box, right=of json] (loader) {\cls{SnapshotLoader}\\\cls{SnapshotValidator}};
    \node[box, right=of loader] (provider) {\cls{SystemStateProvider}};
    \node[box, below=of provider] (filtered) {\cls{FilteringSystemStateProvider}};
    \node[policy, left=of filtered] (prefilter) {\cls{CandidatePrefilter}};

    \node[widebox, right=1.45cm of provider] (manager) {\cls{TemporalWindowManager}};
    \node[box, below=of manager] (dynamicity) {\cls{DynamicityEvaluator}};
    \node[box, below=of dynamicity] (adapter) {\cls{PopulationAdapter}};

    \node[widebox, right=1.45cm of manager] (optimizer) {\cls{MaGaOptimizer}};
    \node[policy, below=of optimizer] (operators) {operatori GA\\policy + repair};
    \node[result, right=1.45cm of optimizer] (result) {\cls{MaGaResult}\\breakdown};

    \draw[arrow] (json) -- (loader);
    \draw[arrow] (loader) -- (provider);
    \draw[arrow] (provider) -- (manager);
    \draw[dashedarrow] (prefilter) -- (filtered);
    \draw[dashedarrow] (provider) -- (filtered);
    \draw[dashedarrow] (filtered) -- (manager);

    \draw[arrow] (manager) -- (dynamicity);
    \draw[arrow] (dynamicity) -- (adapter);
    \draw[arrow] (adapter) -- (optimizer);
    \draw[arrow] (manager) -- (optimizer);
    \draw[arrow] (optimizer) -- (operators);
    \draw[arrow] (optimizer) -- (result);
\end{tikzpicture}
}
\caption{Flusso principale del sistema.}
\end{figure}

Il sistema ha due livelli.

\textbf{Livello GA.} Vive soprattutto nei package \pkg{ga}, \pkg{model} e
\pkg{config}. Lavora su uno snapshot singolo. Produce il miglior cromosoma
trovato.

\textbf{Livello temporale.} Vive nel package \pkg{window}. Decide quando
richiamare il GA, quale snapshot usare, e quanta popolazione della finestra
precedente riusare.

\section{Collegamento con la formalizzazione}

La tabella collega i blocchi matematici al codice. Le frasi sono brevi apposta:
in Overleaf restano leggibili anche su A4.

\begin{table}[h]
\small
\centering
\begin{tabularx}{\linewidth}{@{}L{3.2cm}Y@{}}
\toprule
\textbf{Formalizzazione} & \textbf{Traduzione nel codice} \\
\midrule
Istante o finestra $k$ &
\cls{TemporalWindowManager} esegue uno step. Lo step usa un tempo di trigger e un tempo di osservazione. \\
\addlinespace
Stato $S_k$ &
\cls{SystemSnapshot}. Contiene veicoli, task e candidati disponibili nello stesso istante. \\
\addlinespace
Task $i$ &
\cls{TaskInstance}. Porta id, veicolo sorgente, input, output, cicli CPU e deadline. \\
\addlinespace
Nodo candidato $j$ &
\cls{NodeCandidate}. Descrive nodo locale, veicolo remoto, edge o cloud. Include CPU, banda, latenza e compatibilita' con il sorgente. \\
\addlinespace
Decisione $x_{ij}$ &
Nel codice e' \cls{Gene.selectedCandidateId}. Un gene assegna un task a un candidato. \\
\addlinespace
Quota remota $p_i$ &
\cls{Gene.offloadingRatio}. Se il candidato e' locale vale $0$. Se e' remoto puo' stare tra $0$ e $1$. \\
\addlinespace
Risorse $f_i,b_i$ &
\cls{Gene.allocatedCpu} e \cls{Gene.allocatedBandwidth}. Sono generate e mutate anche da \cls{ResourceAllocationPolicy}. \\
\addlinespace
Soluzione $C$ &
\cls{Chromosome}. E' la lista dei geni, uno per task. \\
\addlinespace
Funzione obiettivo &
\cls{FitnessEvaluator}. Calcola tempo, latenza, mobilita' e penalita' di risorsa. \\
\addlinespace
Evoluzione temporale &
\cls{DynamicityEvaluator} e \cls{PopulationAdapter}. Misurano il cambiamento e decidono il riuso della popolazione. \\
\bottomrule
\end{tabularx}
\caption{Mappa tra formalizzazione e codice.}
\end{table}

\section{Dati del problema}

Le classi del package \pkg{model} non decidono. Rappresentano dati. Questo e'
importante: il GA e le policy devono leggere lo stato, non modificarlo.

\subsection{Snapshot}

\cls{SystemSnapshot} e' il contenitore dello stato. Ha tre liste principali:

\begin{itemize}[leftmargin=*,nosep]
    \item veicoli, tramite \cls{VehicleSnapshot};
    \item task, tramite \cls{TaskInstance};
    \item candidati, tramite \cls{NodeCandidate}.
\end{itemize}

Il GA riceve sempre uno \cls{SystemSnapshot}. Questo evita dipendenze dirette da
JSON, file o provider temporali.

\subsection{Veicoli, task e candidati}

\begin{table}[h]
\small
\centering
\begin{tabularx}{\linewidth}{@{}L{3.4cm}Y@{}}
\toprule
\textbf{Classe} & \textbf{Ruolo} \\
\midrule
\clsb{VehicleSnapshot} &
Descrive un veicolo: posizione, velocita', CPU locale e nodo associato. Serve anche alla stima di copertura. \\
\clsb{TaskInstance} &
Descrive il lavoro da eseguire: sorgente, dimensione input, dimensione output, cicli CPU e deadline. \\
\clsb{NodeCandidate} &
Descrive una possibile destinazione. Puo' essere locale, V2V, edge o cloud. Dice anche se e' valida per il veicolo sorgente. \\
\clsb{CoverageEstimator} &
Stima per quanto tempo il sorgente resta coperto dal candidato. La fitness usa questa stima per la penalita' di mobilita'. \\
\bottomrule
\end{tabularx}
\caption{Dati letti dal GA.}
\end{table}

\section{Cromosoma e gene}

Il cromosoma e' la forma concreta della soluzione.

\[
C = (g_1, g_2, \ldots, g_n)
\]

Ogni gene contiene la decisione per un task:

\[
g_i = (taskId_i, candidateId_i, p_i, f_i, b_i)
\]

\begin{table}[h]
\small
\centering
\begin{tabularx}{\linewidth}{@{}L{3.2cm}Y@{}}
\toprule
\textbf{Campo del gene} & \textbf{Significato} \\
\midrule
\code{taskId} & Task a cui il gene si riferisce. \\
\code{selectedCandidateId} & Nodo scelto per eseguire il task. \\
\code{offloadingRatio} & Quota $p_i$ da eseguire in remoto. Vale $0$ se il candidato e' locale. \\
\code{allocatedCpu} & CPU assegnata al ramo remoto o alla scelta locale. \\
\code{allocatedBandwidth} & Banda assegnata alla trasmissione remota. Vale $0$ nel locale. \\
\bottomrule
\end{tabularx}
\caption{Campi usati dal GA per rappresentare una decisione.}
\end{table}

Questa struttura corrisponde alla formalizzazione: scelta del nodo,
partizionamento del task e allocazione di risorse.

\section{Caricamento e provider}

\cls{SnapshotLoader} legge gli snapshot dai file. \cls{SnapshotValidator}
controlla che lo scenario sia coerente prima di darlo al resto del sistema.

Il gestore temporale non conosce i file. Dipende solo da \cls{SystemStateProvider}.
Questa interfaccia espone un solo metodo:

\begin{lstlisting}[language=Java]
Optional<SystemSnapshot> findSnapshotAt(double observationTimeSeconds);
\end{lstlisting}

\subsection{Correzione del provider filtrato}

\cls{FilteringSystemStateProvider} e' un decorator. Riceve un provider reale e
un \cls{CandidatePrefilter}. Quando il manager chiede uno snapshot, il decorator:

\begin{enumerate}[leftmargin=*,nosep]
    \item chiama il provider delegato con \meth{findSnapshotAt};
    \item applica il prefilter allo snapshot trovato;
    \item salva il risultato diagnostico;
    \item restituisce lo snapshot filtrato.
\end{enumerate}

La correzione recente e' questa: il provider filtrato implementa solo
\meth{findSnapshotAt(double)}. Non espone piu' un override di
\meth{findSnapshotAtOrAfter}, perche' quel metodo non fa parte del contratto
\cls{SystemStateProvider}. Cosi' il decorator resta allineato
all'interfaccia e non chiede al delegato un metodo che non deve avere.

\section{Prefilter dei candidati}

Il prefilter non sceglie la soluzione. Riduce lo spazio di ricerca prima del GA.
Lo fa per evitare che il GA sprechi generazioni su candidati chiaramente deboli.

\begin{figure}[h]
\centering
\resizebox{0.80\linewidth}{!}{%
\begin{tikzpicture}[node distance=1.0cm and 1.1cm]
    \node[data] (candidate) {candidato remoto};
    \node[decision, right=of candidate] (source) {ha task\\compatibili?};
    \node[decision, right=of source] (resource) {CPU e banda\\sufficienti?};
    \node[decision, right=of resource] (coverage) {copertura\\sufficiente?};
    \node[decision, right=of coverage] (deadline) {lower bound\\entro deadline?};
    \node[result, right=of deadline] (keep) {mantieni};
    \node[result, below=of resource] (drop) {scarta\\o fallback};

    \draw[arrow] (candidate) -- (source);
    \draw[arrow] (source) -- node[above,font=\scriptsize]{si} (resource);
    \draw[arrow] (resource) -- node[above,font=\scriptsize]{si} (coverage);
    \draw[arrow] (coverage) -- node[above,font=\scriptsize]{si} (deadline);
    \draw[arrow] (deadline) -- node[above,font=\scriptsize]{si} (keep);
    \draw[arrow] (source) |- node[left,font=\scriptsize]{no} (drop);
    \draw[arrow] (resource) -- node[right,font=\scriptsize]{no} (drop);
    \draw[arrow] (coverage) |- node[right,font=\scriptsize]{no} (drop);
    \draw[arrow] (deadline) |- node[right,font=\scriptsize]{no} (drop);
\end{tikzpicture}
}
\caption{Decisione del prefilter.}
\end{figure}

Il prefilter produce anche un report. Le ragioni principali sono:
\code{NO_TASK_FOR_SOURCE}, \code{INVALID_CPU}, \code{INVALID_BANDWIDTH},
\code{INSUFFICIENT_COVERAGE}, \code{DEADLINE_LOWER_BOUND_TOO_HIGH} e
\code{RESTORED_AS_FALLBACK}. L'ultimo caso e' una protezione: un task non deve
restare senza alcun candidato.

\section{Policy di offloading}

\cls{OffloadingRatioPolicy} decide valori plausibili per $p_i$. Non sceglie il
nodo. Non valuta la soluzione. Aiuta solo inizializzazione e mutazione a
generare quote migliori.

La policy usa piu' modalita':

\begin{itemize}[leftmargin=*,nosep]
    \item quota casuale, per mantenere esplorazione;
    \item quota bilanciata, per avvicinare ramo locale e ramo remoto;
    \item quota deadline-aware, per tenere conto della deadline;
    \item quota pienamente remota, utile come prova su candidati promettenti;
    \item piccolo passo locale, utile in mutazione.
\end{itemize}

La parte nuova e' \meth{deadlineAwareRatio}. Questa stima se il remoto puo'
aiutare davvero. Considera il peso dell'upload e usa una soglia di dominanza
dell'upload. Se il remoto e' troppo lento gia' nel lower bound, torna a una
quota casuale. Quindi non forza il full offloading quando non ha senso.

\section{Policy di CPU e banda}

Le classi nuove sono \cls{ResourceAllocationPolicy} e
\cls{ResourceAllocationDecision}. Sono centrali per capire le correzioni recenti.

\cls{ResourceAllocationPolicy} genera o muta CPU e banda per un gene remoto.
Guarda quattro cose:

\begin{itemize}[leftmargin=*,nosep]
    \item deadline del task;
    \item quota $p_i$;
    \item capacita' del candidato;
    \item stima minima dei tempi locale e remoto.
\end{itemize}

Poi classifica la scelta come fattibile, borderline o non fattibile. Questa
classificazione non sostituisce la fitness. Serve solo a evitare allocazioni
assurde, per esempio saturare CPU e banda su un candidato che non puo'
rispettare la deadline nemmeno in condizioni ottimistiche.

\begin{figure}[h]
\centering
\resizebox{0.86\linewidth}{!}{%
\begin{tikzpicture}[node distance=1.0cm and 1.2cm]
    \node[data] (input) {task\\candidato\\veicolo\\$p_i$};
    \node[policy, right=of input] (policy) {\cls{ResourceAllocationPolicy}};
    \node[decision, right=of policy] (feas) {fattibilita'};
    \node[policy, right=of feas] (mode) {modo\\di allocazione};
    \node[result, right=of mode] (decision) {\cls{ResourceAllocationDecision}\\CPU, banda, modo};

    \draw[arrow] (input) -- (policy);
    \draw[arrow] (policy) -- (feas);
    \draw[arrow] (feas) -- (mode);
    \draw[arrow] (mode) -- (decision);
\end{tikzpicture}
}
\caption{Da vincoli locali a decisione di risorsa.}
\end{figure}

\cls{ResourceAllocationDecision} e' immutabile. Contiene:
\code{allocatedCpu}, \code{allocatedBandwidth} e \code{Mode}. I modi sono
\code{LOCAL}, \code{RANDOM}, \code{DEADLINE_AWARE}, \code{BORDERLINE},
\code{MODERATE}, \code{AGGRESSIVE} e \code{SMALL_STEP}.

Questa informazione e' utile anche per la diagnostica. Se molti geni finiscono
in \code{MODERATE} o \code{BORDERLINE}, lo scenario e' probabilmente stretto.

\section{Inizializzazione e mutazione}

\subsection{PopulationInitializer}

\cls{PopulationInitializer} crea la popolazione iniziale. Ora non usa un solo
schema casuale. Alterna quattro profili:

\begin{table}[h]
\small
\centering
\begin{tabularx}{0.92\linewidth}{@{}L{3.1cm}Y@{}}
\toprule
\textbf{Profilo} & \textbf{Idea} \\
\midrule
\code{LOCAL_BIASED} & Circa 15\%. Tiene molte scelte locali, utile come base robusta. \\
\code{BALANCED_REMOTE} & Circa 35\%. Usa quote deadline-aware su candidati remoti validi. \\
\code{FULL_REMOTE_TRIAL} & Circa 10\%. Prova $p_i=1$ su candidati remoti. \\
\code{RANDOM} & Circa 40\%. Mantiene diversita' genetica. \\
\bottomrule
\end{tabularx}
\caption{Profili usati per costruire la popolazione iniziale.}
\end{table}

Per i geni remoti, CPU e banda passano da
\meth{ResourceAllocationPolicy.allocateInitial}. Quindi l'inizializzazione non
assegna piu' risorse solo con un numero casuale.

\subsection{MutationOperator}

\cls{MutationOperator} modifica i geni durante il GA. La mutazione resta
casuale, ma ora e' piu' guidata.

Quando cambia candidato, preferisce un candidato remoto con probabilita' $0.60$.
Tra i remoti puo' scegliere quello con migliore stima euristica con probabilita'
$0.55$. Quando cambia $p_i$, puo' usare \meth{deadlineAwareRatio}. Quando cambia
CPU e banda, usa \meth{ResourceAllocationPolicy.mutate}.

Il punto e' semplice: la mutazione non deve solo fare rumore. Deve anche provare
mosse che hanno una possibilita' concreta di migliorare la fitness.

\section{Repair}

\cls{RepairOperator} controlla che ogni gene sia coerente con lo snapshot.
Corregge:

\begin{itemize}[leftmargin=*,nosep]
    \item task mancanti nel cromosoma;
    \item candidato inesistente o non valido per il veicolo sorgente;
    \item quota fuori dall'intervallo $[0,1]$;
    \item CPU e banda negative o oltre la capacita' del singolo candidato;
    \item gene locale, che deve avere $p_i=0$ e banda $0$.
\end{itemize}

Dopo il repair del singolo gene entra \cls{CpuAggregateRepairOperator}. Questa
classe ridimensiona la CPU aggregata per nodo fisico remoto. Serve per evitare
che piu' candidati legati allo stesso nodo sommino piu' CPU di quella
disponibile.

La banda aggregata resta una parte aperta: il codice la limita sul singolo
candidato e la penalizza in fitness, ma non ha ancora un repair aggregato
simmetrico alla CPU.

\section{Fitness}

\cls{FitnessEvaluator} traduce un cromosoma in un numero. Piu' il numero e'
basso, migliore e' la soluzione.

La struttura segue la formalizzazione:

\[
J(C) =
w_T \hat{T}(C) +
w_L \hat{L}(C) +
w_M \hat{M}(C) +
w_R \hat{R}(C)
\]

Dove:

\begin{itemize}[leftmargin=*,nosep]
    \item $\hat{T}(C)$ e' il tempo di completamento normalizzato;
    \item $\hat{L}(C)$ e' la latenza media di comunicazione normalizzata;
    \item $\hat{M}(C)$ e' la penalita' di mobilita' normalizzata;
    \item $\hat{R}(C)$ include penalita' di risorsa, vincoli e soluzioni invalide.
\end{itemize}

\subsection{Tempo di un gene}

Per un gene locale:

\[
T_i = \frac{cycles_i}{f_i^{local}}
\]

Per un gene remoto:

\[
T_i^{remote} =
T_i^{upload} + T_i^{exec} + T_i^{download} + L_i^{base}
\]

Se il task e' parziale, il codice considera anche il ramo locale. Il tempo del
gene e' il massimo tra ramo locale e ramo remoto:

\[
T_i = \max(T_i^{local}, T_i^{remote})
\]

Il tempo del cromosoma e' il massimo sui task:

\[
T(C)=\max_i T_i
\]

\subsection{Latenza e mobilita'}

La latenza di comunicazione remota include upload, download e latenza base del
candidato:

\[
L_i = T_i^{upload} + T_i^{download} + L_i^{base}
\]

Nel codice la latenza usata dalla fitness e' la media sui task, non la somma.
Quindi:

\[
L(C)=\frac{1}{|T|}\sum_i L_i
\]

La mobilita' usa due parti attive: rischio di copertura e rischio di handover.
Nel codice attuale \code{linkInstability} vale $0.0$. E' quindi un segnaposto,
non una componente misurata.

\subsection{Risorse}

La CPU remota viene aggregata per \code{executionNodeId}. La banda viene
aggregata per \code{candidateId}. Questo dettaglio conta: piu' candidati possono
riferirsi allo stesso nodo fisico per la CPU, mentre la banda resta legata al
link candidato.

\section{Ciclo del GA}

\cls{MaGaOptimizer} coordina il ciclo genetico. Riceve uno snapshot e, se
disponibile, una popolazione iniziale gia' adattata dal livello temporale.

\begin{figure}[h]
\centering
\resizebox{0.78\linewidth}{!}{%
\begin{tikzpicture}[node distance=1.0cm and 1.2cm]
    \node[data] (start) {snapshot\\popolazione iniziale};
    \node[box, below=of start] (repair0) {repair\\fitness};
    \node[box, below=of repair0] (select) {selection};
    \node[box, below=of select] (cross) {crossover};
    \node[box, below=of cross] (mut) {mutation};
    \node[box, below=of mut] (repair) {repair\\fitness};
    \node[decision, below=of repair] (stop) {stop?};
    \node[result, below=of stop] (out) {miglior cromosoma\\popolazione finale};

    \draw[arrow] (start) -- (repair0);
    \draw[arrow] (repair0) -- (select);
    \draw[arrow] (select) -- (cross);
    \draw[arrow] (cross) -- (mut);
    \draw[arrow] (mut) -- (repair);
    \draw[arrow] (repair) -- (stop);
    \draw[arrow] (stop) -- node[right,font=\scriptsize]{si} (out);
    \draw[arrow] (stop.west) -- ++(-1.6,0) |- node[left,font=\scriptsize]{no} (select.west);
\end{tikzpicture}
}
\caption{Ciclo genetico su uno snapshot.}
\end{figure}

Gli operatori principali sono:
\cls{SelectionOperator}, \cls{CrossoverOperator}, \cls{MutationOperator},
\cls{RepairOperator} ed \cls{ElitismOperator}. Il ciclo si ferma per numero
massimo di generazioni o stagnazione. Il risultato conserva anche statistiche
di generazione e breakdown della fitness.

\section{Gestione temporale}

\cls{TemporalWindowManager} avvia il GA piu' volte lungo la linea temporale.
Il flusso e':

\[
trigger \rightarrow observationTime \rightarrow snapshot
\rightarrow dynamicity \rightarrow reuseMode \rightarrow popolazione
\rightarrow MA\text{-}GA
\]

Il tempo osservato e':

\[
observationTime = triggerTime + dataCollectionDelay
\]

Il trigger puo' avere tre motivi:
\code{FIRST_RUN}, \code{SCHEDULED_WINDOW_EXPIRATION} o \code{CRITICAL_EVENT}.

Nel codice attuale l'intervallo programmato usa
\meth{windowConfig.getFixedIntervalSeconds}. La dinamicita' non cambia ancora la
durata della finestra. Cambia invece il riuso della popolazione.

\section{Eventi critici}

\cls{CriticalEvent} rappresenta un fatto che puo' anticipare la prossima
ottimizzazione. Non e' uno snapshot. E' un segnale: qualcosa e' cambiato e la
soluzione corrente potrebbe non essere piu' buona.

Gli eventi hanno:

\begin{itemize}[leftmargin=*,nosep]
    \item un tempo;
    \item un tipo, per esempio link perso, task arrivato o rischio deadline;
    \item una severita';
    \item un riferimento opzionale a veicolo, task, candidato o nodo.
\end{itemize}

\cls{CriticalEventSeverity} decide se l'evento deve attivare la
riottimizzazione. \code{HIGH} e \code{CRITICAL} la attivano. \code{LOW} e
\code{MEDIUM} restano informativi.

\section{Dinamicita' e riuso}

\cls{DynamicityEvaluator} confronta lo snapshot precedente con quello corrente.
Produce un \cls{DynamicityBreakdown}. Il breakdown contiene il valore globale,
i contributi per veicoli, task, risorse e link, e il livello qualitativo.

\begin{table}[h]
\small
\centering
\begin{tabularx}{0.92\linewidth}{@{}L{2.8cm}Y@{}}
\toprule
\textbf{Livello} & \textbf{Modo di riuso} \\
\midrule
\code{UNKNOWN} & \code{FIRST_RUN}. Non esiste uno snapshot precedente. \\
\code{STABLE} & \code{WARM_START}. Si prova a riusare quasi tutta la popolazione precedente. \\
\code{MODERATE} & \code{PARTIAL_RESTART}. Si mantiene una quota dei migliori cromosomi e si rigenera il resto. \\
\code{HIGH} & \code{COLD_START}. Lo scenario e' cambiato troppo. Si riparte da popolazione nuova. \\
\bottomrule
\end{tabularx}
\caption{Mappatura tra dinamicita' e riuso.}
\end{table}

\cls{PopulationAdapter} applica questa decisione. Copia i cromosomi precedenti,
li adatta allo snapshot corrente, li ripara e completa la popolazione quando
serve.

\section{Catene di metodi principali}

Questa sezione collega le classi ai metodi. E' utile quando segui il codice nel
debugger o dentro l'IDE.

\begin{table}[h]
\small
\centering
\begin{tabularx}{\linewidth}{@{}L{3.2cm}Y@{}}
\toprule
\textbf{Percorso} & \textbf{Chiamate e motivo} \\
\midrule
Finestra temporale &
\meth{run} chiama \meth{executeNextStepOrNull}. Dentro lo step compaiono \meth{resolveTrigger}, \meth{findSnapshotAt}, \meth{evaluate}, \meth{adaptPopulation} e \meth{optimizeDetailed}. Questo percorso decide quando ottimizzare e con quale popolazione partire. \\
\addlinespace
Provider filtrato &
\meth{findSnapshotAt} chiama il provider delegato, poi \meth{filterSnapshot}. Il filtro usa \meth{CandidatePrefilter.filter} e restituisce \meth{getFilteredSnapshot}. Questo percorso inserisce il prefilter senza cambiare il manager. \\
\addlinespace
Inizializzazione &
\meth{PopulationInitializer} crea un gene scegliendo un profilo. Per i remoti usa \meth{deadlineAwareRatio} e \meth{allocateInitial}. Questo percorso prepara cromosomi meno casuali. \\
\addlinespace
Mutazione &
\meth{MutationOperator.mutate} visita i geni. Se serve cambia candidato, quota e risorse. Per la quota usa \meth{deadlineAwareRatio}; per CPU e banda usa \meth{ResourceAllocationPolicy.mutate}. \\
\addlinespace
Valutazione &
\meth{evaluateDetailed} scorre i task, chiama la valutazione del gene e poi aggrega tempo, latenza, mobilita' e risorse. Questo percorso produce \cls{EvaluationBreakdown}. \\
\addlinespace
Diagnostica &
\cls{DeadlineViolationAnalyzer} legge i breakdown dopo la fitness. Non cambia il risultato: spiega perche' una deadline e' stata violata. \\
\bottomrule
\end{tabularx}
\caption{Collegamento tra metodi e responsabilita'.}
\end{table}

\section{Risultati e diagnostica}

Il sistema non restituisce solo il miglior valore di fitness. Produce oggetti
che spiegano il perche' del risultato.

\begin{table}[h]
\small
\centering
\begin{tabularx}{\linewidth}{@{}L{3.5cm}Y@{}}
\toprule
\textbf{Classe} & \textbf{Cosa spiega} \\
\midrule
\clsb{MaGaResult} &
Miglior cromosoma, fitness migliore, stop reason, statistiche e popolazione finale. \\
\clsb{EvaluationBreakdown} &
Scomposizione della fitness: tempo, latenza, mobilita', risorse e penalita'. \\
\clsb{GeneEvaluationBreakdown} &
Dettaglio per task: tempi locale/remoto, latenza, deadline e penalita'. \\
\clsb{CandidateFilteringResult} &
Snapshot filtrato e motivi di rimozione dei candidati. \\
\clsb{DeadlineViolationAnalyzer} &
Diagnostica post-fitness sulle violazioni di deadline. Non modifica cromosomi. \\
\clsb{TemporalStepResult} &
Risultato di una finestra temporale: trigger, snapshot, dinamicita', riuso e risultato GA. \\
\bottomrule
\end{tabularx}
\caption{Oggetti usati per capire il risultato.}
\end{table}

\cls{DeadlineViolationAnalyzer} serve dopo l'ottimizzazione. Legge i breakdown
e classifica la causa del problema: collo di bottiglia locale, upload, remoto,
download, latenza base, copertura o caso misto. Questo aiuta a capire se il
problema viene dai dati, dalle risorse o dalla mobilita'.

\section{Cosa decide ogni classe}

Questa tabella riassume le responsabilita'. E' la parte da usare quando vuoi
capire ``chi fa cosa''.

\begin{table}[h]
\small
\centering
\begin{tabularx}{\linewidth}{@{}L{3.4cm}Y@{}}
\toprule
\textbf{Classe} & \textbf{Decisione o responsabilita'} \\
\midrule
\clsb{SystemStateProvider} &
Fornisce lo snapshot osservato a un tempo simulato. \\
\clsb{FilteringSystemStateProvider} &
Applica il prefilter prima di restituire lo snapshot al manager. \\
\clsb{CandidatePrefilter} &
Rimuove candidati remoti con risorse, copertura o lower bound troppo deboli. \\
\clsb{TemporalWindowManager} &
Decide quando rieseguire il GA e quale snapshot usare. \\
\clsb{DynamicityEvaluator} &
Misura quanto e' cambiato lo scenario. \\
\clsb{PopulationAdapter} &
Decide come riusare o rigenerare la popolazione precedente. \\
\clsb{PopulationInitializer} &
Crea nuovi cromosomi con profili locali, remoti bilanciati, full remote e random. \\
\clsb{OffloadingRatioPolicy} &
Propone quote $p_i$ per inizializzazione e mutazione. \\
\clsb{ResourceAllocationPolicy} &
Propone CPU e banda coerenti con deadline, quota e capacita'. \\
\clsb{MutationOperator} &
Introduce variazioni nei geni, usando anche le policy nuove. \\
\clsb{RepairOperator} &
Riporta cromosomi e geni dentro i vincoli base. \\
\clsb{FitnessEvaluator} &
Valuta la soluzione e decide quanto costa. \\
\clsb{MaGaOptimizer} &
Coordina il ciclo genetico e restituisce il miglior risultato. \\
\bottomrule
\end{tabularx}
\caption{Responsabilita' principali.}
\end{table}

\section{Differenze operative rispetto alla formalizzazione}

Il codice segue la struttura della formalizzazione, ma alcune scelte sono piu'
operative. Conviene tenerle presenti.

\begin{table}[h]
\small
\centering
\begin{tabularx}{\linewidth}{@{}L{3.6cm}Y@{}}
\toprule
\textbf{Punto} & \textbf{Stato nel codice} \\
\midrule
Latenza &
La formalizzazione puo' leggerla come somma. Il codice usa la media sui task. \\
\addlinespace
Instabilita' link &
La variabile esiste nella formula della mobilita', ma nel codice attuale vale $0.0$. \\
\addlinespace
Intervallo temporale &
Il manager usa un intervallo fisso. La dinamicita' influisce sul riuso della popolazione, non sulla durata della finestra. \\
\addlinespace
Prefilter &
Nel codice e' gia' integrato tramite \cls{FilteringSystemStateProvider}. Riduce candidati prima del GA. \\
\addlinespace
Quote e risorse &
\cls{OffloadingRatioPolicy} e \cls{ResourceAllocationPolicy} guidano il campionamento. La fitness resta il giudice finale. \\
\addlinespace
Banda aggregata &
La fitness la penalizza. Il repair aggregato simmetrico alla CPU non e' ancora implementato. \\
\bottomrule
\end{tabularx}
\caption{Punti in cui il codice va letto con attenzione.}
\end{table}

\section{Percorso consigliato di lettura del codice}

Per capire il progetto senza perdersi, conviene seguire questo ordine:

\begin{enumerate}[leftmargin=*]
    \item \cls{SystemSnapshot}, \cls{TaskInstance}, \cls{VehicleSnapshot}, \cls{NodeCandidate}.
    \item \cls{Gene} e \cls{Chromosome}.
    \item \cls{FitnessEvaluator}, leggendo anche \cls{EvaluationBreakdown}.
    \item \cls{OffloadingRatioPolicy} e \cls{ResourceAllocationPolicy}.
    \item \cls{PopulationInitializer}, \cls{MutationOperator} e \cls{RepairOperator}.
    \item \cls{MaGaOptimizer}.
    \item \cls{TemporalWindowManager}, \cls{DynamicityEvaluator} e \cls{PopulationAdapter}.
    \item \cls{CandidatePrefilter} e \cls{DeadlineViolationAnalyzer}.
\end{enumerate}

Questo ordine parte dai dati, passa per una singola soluzione, poi arriva al GA
e solo alla fine al ciclo temporale.

\section{Sintesi finale}

Il codice realizza la formalizzazione in modo modulare. Lo snapshot e' lo stato.
Il gene e' la decisione elementare. Il cromosoma e' la strategia completa. Le
policy propongono scelte iniziali e mutazioni piu' sensate. Il repair tiene le
soluzioni dentro i vincoli base. La fitness misura il costo. Il GA cerca il
cromosoma migliore. Il livello temporale decide quando ripetere il processo e
quanto riusare della popolazione precedente.

Le correzioni recenti rendono il sistema piu' coerente. Il provider filtrato ora
rispetta il contratto di \cls{SystemStateProvider}. Le nuove policy di quota e
risorse evitano molte scelte casuali deboli senza togliere esplorazione al GA.
La decisione finale resta comunque nella fitness e nella selezione genetica.

\end{document}
